\PassOptionsToPackage{hidelinks}{hyperref}
\documentclass[a4paper,10pt]{oblivoir}
\usepackage{amsmath,amssymb,amsfonts}
\usepackage{graphicx}
\usepackage{booktabs}
\usepackage{array}
\usepackage{setspace}
\usepackage{geometry}
\usepackage{caption}
\usepackage{float}
\geometry{a4paper, top=25mm, bottom=25mm, left=25mm, right=25mm}
\setstretch{1.8}
\setcounter{secnumdepth}{0}
\graphicspath{{./}}

\begin{document}
\title{중국인 유학생의 한국어 쓰기 태도 인과 구조: 베이지안 부분 신용 모형–구조 방정식 모형(PCM-SEM) 분석}
\maketitle

\textbf{Causal Structure of Korean Writing Attitudes among Chinese International Students: A Bayesian Partial Credit Model–Structural Equation Model (PCM-SEM) Analysis}


\subsection{초록}

[데이터 수집 및 분석 완료 후 작성]

\textbf{주제어}: 중국인 유학생, 한국어 쓰기 태도, 베이지안 PCM-SEM, 인과 매개, 완전 매개, MCMC


\subsection{1. 서론}

한국 내 중국인 유학생 수는 2023년 기준 약 67,000명으로 전체 외국인 유학생의 44\%를 차지한다(교육부, 2023). 이들 대부분은 학문 목적으로 한국어를 학습하며, 한국어 쓰기 능력은 대학 학업 수행의 핵심 역량이다. 그러나 중국인 유학생의 쓰기 태도를 인과 경로의 관점에서 체계적으로 분석한 연구는 드물다.

주월랑(2022)은 외국인 유학생 86명(13개국 출신)을 대상으로 쓰기 태도를 쓰기인식(Cognition), 쓰기반응(Affect), 수행태도(Behavior)의 세 하위 구인으로 측정하고 학습자 배경 변수에 따른 집단 차이를 분석하였다. 그러나 이 연구는 세 구인 간의 인과 경로를 직접 검증하지 않았으며, 방법론적으로 합산 점수 OLS를 사용함으로써 문항 수준 측정 오차를 무시하는 회귀 감쇠 편향의 문제를 안고 있다.

본 연구는 두 가지 측면에서 선행 연구를 확장한다. 첫째, 연구 대상을 중국인 유학생으로 집중하여 이 집단의 쓰기 태도 인과 구조를 규명한다. 중국인 유학생은 한국어 쓰기 교육의 가장 중요한 수혜 집단임에도 불구하고, 이들의 쓰기 태도 메커니즘에 대한 독립적 연구는 부족하다. 둘째, 별도의 방법론 연구(연구 1; t1\_paper 참조)에서 타당화된 베이지안 부분 신용 모형–구조 방정식 모형(PCM-SEM)을 적용하여, 문항 수준 측정 오차를 통제하고 인과 경로의 방향성과 크기를 확률적으로 추정한다.

연구 문제는 다음과 같다.

RQ1. 중국인 유학생의 쓰기인식이 쓰기반응에 미치는 영향($\beta_1$)은 유의한가?

RQ2. 쓰기반응이 수행태도에 미치는 영향($\beta_2$)은 유의한가?

RQ3. 쓰기인식이 수행태도에 미치는 직접 효과($\gamma_1$)는 쓰기반응의 완전 매개 후 유의하지 않게 되는가?

RQ4. 성별이 쓰기반응 및 수행태도에 미치는 영향은 유의한가?


\subsection{2. 이론적 배경}

\subsubsection{2.1 쓰기 태도의 구조와 인과 경로}

본 연구는 쓰기 태도를 인지적·정서적·행동적 요소가 위계적으로 연결된 구조로 상정한다. 이를 위해 태도 구성 요소의 ABC 모델(Ajzen \& Fishbein, 1977)과 제2언어 습득 이론의 정의적 여과기(Affective Filter) 가설(Krashen, 1982)을 결합하여 다음과 같은 인과 경로를 설정하였다.

\textbf{쓰기인식(Cognition, $X$)}은 쓰기의 학문적 가치와 필요성에 대한 지적 판단을 의미하며, 태도 형성의 출발점이다. \textbf{쓰기반응(Affect, $M$)}은 쓰기에 대한 즐거움, 자신감, 불안 등 정서적 차원을 포함한다. 인지적 가치 인식이 정서적 반응을 매개하지 않으면 실제 행동 변화로 이어지기 어렵다. \textbf{수행태도(Behavior, $Y$)}는 독자 고려, 고쳐쓰기 등 실제 쓰기 수행에 관한 행동적 의지를 나타낸다.

\[X(\text{쓰기인식}) \rightarrow M(\text{쓰기반응}) \rightarrow Y(\text{수행태도})\]

자기결정성 이론(Self-Determination Theory; Deci \& Ryan, 1985)의 관점에서도 내재적 동기 형성에 정의적 반응이 필수적 매개 역할을 담당한다. 쓰기의 중요성을 인지하더라도 정서적으로 긍정적 반응을 형성하지 못하면 능동적 수행태도로 전환되지 않는다는 매개 경로 가설이 수립된다.

\subsubsection{2.2 중국인 유학생과 한국어 쓰기}

중국인 유학생의 한국어 쓰기 학습 맥락은 몇 가지 고유한 특징을 가진다. 중국어와 한국어는 표기 체계가 상이하며, 학문적 글쓰기의 문화적 관습—서론-본론-결론 구조의 명시성, 인용 방식 등—에서 차이가 있다(Kaplan, 1966). 또한 중국인 유학생의 한국어 쓰기 동기는 학업 요건 충족이라는 외재적 동기가 강하게 작용할 수 있으며, 이는 쓰기반응(정서적 차원)과 수행태도(행동 의지) 간의 연결 구조에 영향을 미칠 수 있다.

이러한 맥락을 고려할 때, 중국인 유학생 집단에서도 쓰기인식\ensuremath{\rightarrow}쓰기반응\ensuremath{\rightarrow}수행태도의 순차적 인과 경로가 작동하는지를 검증하는 것은 교육적으로 중요한 의미를 가진다.

\subsubsection{2.3 베이지안 PCM-SEM의 역할}

합산 점수 OLS는 측정 오차로 인한 회귀 감쇠 편향을 내포하며, 이 편향은 표본 크기가 증가해도 해소되지 않는다(연구 1 참조). PCM-SEM은 순서형 문항 반응을 부분 신용 모형(PCM; Masters, 1982)으로 직접 모형화함으로써 이 문제를 해결한다.

베이지안 추론은 사전 분포를 통해 선행 연구 지식을 체계적으로 통합하며, 사후 분포에서 직접 확률적 진술이 가능하다. 소표본(N\ensuremath{\approx}100) 조건에서 사전 정보의 활용은 추정 정밀도를 실질적으로 향상시킨다(Bernardo \& Smith, 1994).

\subsubsection{2.4 연구 가설}

선행 이론과 주월랑(2022)의 기술통계를 바탕으로 다음 가설을 설정한다.

\textbf{H1}: 쓰기인식($X$)은 쓰기반응($M$)에 유의한 정적 영향을 미친다 ($\beta_1 > 0$).

\textbf{H2}: 쓰기반응($M$)은 수행태도($Y$)에 유의한 정적 영향을 미친다 ($\beta_2 > 0$).

\textbf{H3}: 쓰기인식이 수행태도에 미치는 직접 효과($\gamma_1$)는 유의하지 않다 (완전 매개).

\textbf{H4}: 쓰기인식이 쓰기반응을 거쳐 수행태도에 미치는 간접 효과($\beta_1\beta_2$)는 유의하다.

\textbf{H5}: 성별이 쓰기반응 및 수행태도에 미치는 영향은 유의하지 않다.


\subsection{3. 연구 방법}

\subsubsection{3.1 연구 참여자}

연구 참여자는 한국 소재 대학교에 재학 중인 중국 국적 유학생 약 100명이다. 선정 기준은 (1) 중국 국적 보유, (2) 한국어 관련 수업 수강 중, (3) 연구 참여 자발적 동의이다. 표본 크기 N\ensuremath{\approx}100은 연구 1의 몬테카를로 분석에서 PCM-SEM이 안정적인 추정 성능을 보인 조건(N=100, 검출력 \ensuremath{\approx} 0.88\~{}0.94)에 해당한다.

인구통계학적 특성(성별, 학년, 한국어 학습 기간, 전공)은 설문지 전면부에서 수집한다.

\subsubsection{3.2 측정 도구}

\textbf{한국어 쓰기 태도 척도}: 주월랑(2022)의 21개 문항(5점 Likert)을 원형 그대로 사용한다. 문항 구성은 쓰기인식 4문항(y1–y4), 쓰기반응 11문항(y5–y15), 수행태도 6문항(y16–y21)이며, 원논문에서 전체 Cronbach $\alpha = 0.854$로 보고되었다.

문항은 "한국어로 글을 쓰는 것이 중요하다고 생각한다"(쓰기인식), "한국어로 글을 쓸 때 즐겁다"(쓰기반응), "한국어 글쓰기에서 독자를 고려한다"(수행태도)와 같은 형식이며, 원논문의 검증된 문항을 그대로 적용함으로써 측정 동등성을 확보한다.

\textbf{응답 형식}: 1점(전혀 그렇지 않다) \~{} 5점(매우 그렇다).

\subsubsection{3.3 분석 방법}

\textbf{모형 보정}: 연구 1(t1 연구)의 몬테카를로 및 시뮬레이션 결과를 바탕으로 PCM 임계값 오프셋 초기값을 설정한다($c_X=-1.194$, $c_M=-0.380$, $c_Y=-0.185$; \texttt{t2\\_calibration.json} 참조). 이 보정값은 중국인 유학생 표본에서도 원논문 수준의 응답 분포를 모형이 예측할 수 있는지를 사전 예측 분포 점검(\texttt{t2\\_model\\_check.py})을 통해 검증한 후 사용한다.

\textbf{Stan 모형}: 연구 1에서 타당화된 두 Stan 모형—약한 사전 분포(\texttt{sem\\_pcm\\_v2.stan})와 강한 사전 분포(\texttt{sem\\_pcm\\_with\\_prior.stan})—을 그대로 적용한다. 강한 사전 분포는 원논문의 구인별 평균과 문헌 효과 크기($r \approx 0.30$\~{}$0.55$)를 반영한다.

\textbf{MCMC 설정}: 4개 체인, 워밍업 1,000회, 샘플링 1,000회 (총 4,000 유효 샘플), adapt\_delta=0.92, max\_treedepth=12.

\textbf{수렴 진단}: $\hat{R} \leq 1.01$, ESS\_bulk $\geq 400$, 발산 전이 0건을 수렴의 기준으로 한다.

\textbf{윤리}: 본 연구는 소속 기관 생명윤리위원회(IRB) 심의를 거쳐 수행한다. 참여자에게 연구 목적과 익명성 보장을 설명하고 서면 동의를 받는다.


\subsection{4. 결과}

[데이터 수집 및 MCMC 분석(t2\_analysis.py) 완료 후 작성]

아래 소절 구조에 따라 결과를 기술한다.

\subsubsection{4.1 참여자 기술통계}

\begin{itemize}
  \item 전체 N, 성별 분포, 학년, 한국어 학습 기간
  \item 구인별 복합 점수 평균 및 표준편차
  \item Cronbach α (구인별)
\end{itemize}

\subsubsection{4.2 MCMC 수렴 진단}

\begin{itemize}
  \item $\hat{R}$, ESS\_bulk, 발산 전이 수 (t2\_convergence.csv 기반)
  \item 표 형식 제시
\end{itemize}

\subsubsection{4.3 구조 경로 계수}

\begin{itemize}
  \item 표: 약한 vs 강한 사전 분포 모형 비교 (사후 평균, SD, 95\% CrI, P(>0))
  \item RQ1–RQ4에 대한 검증 결과
  \item 그림: t2\_fig\_posterior.png, t2\_fig\_path\_estimates.png
\end{itemize}

\subsubsection{4.4 매개 효과 분석}

\begin{itemize}
  \item 간접 효과 $\beta_1\beta_2$ 사후 분포
  \item 완전 매개 vs 부분 매개 판단
  \item $P(\beta_1>0 \wedge \beta_2>0 \mid \text{data})$
  \item 그림: t2\_fig\_mediation.png
\end{itemize}


\subsection{5. 논의}

[데이터 수집 및 분석 완료 후 작성]

아래 주제를 중심으로 논의한다.

\subsubsection{5.1 연구 가설 검증 해석}

H1–H5 각각에 대한 결과 해석 및 이론적 함의. 완전 매개 결과가 확인될 경우: "쓰기의 가치를 인식하더라도 쓰기에 대한 긍정적 정서가 형성되지 않으면 능동적 수행태도로 연결되지 않는다"는 ABC 모델 및 Affective Filter 이론의 경험적 지지.

\subsubsection{5.2 선행 연구(주월랑, 2022)와의 관계}

연구 1(시뮬레이션)과 본 연구(실측) 결과의 비교 해석. 동일 측정 도구를 사용했을 때 중국인 유학생 표본에서 원논문과 유사한 인과 구조가 재현되는지 여부.

\subsubsection{5.3 사전 분포와 소표본 추론}

약한 vs 강한 사전 분포 간 결론 차이가 있다면 그 의미. 원논문 통계 기반 사전 분포의 정보적 가치.

\subsubsection{5.4 교육적 함의}

중국인 유학생 대상 한국어 쓰기 교육 개선을 위한 구체적 제언. 완전 매개가 확인될 경우 인지적 개입(쓰기 가치 인식 제고)만으로는 불충분하며, 정의적 개입(쓰기 불안 감소, 자신감 강화 활동)이 병행되어야 함.

\subsubsection{5.5 연구의 한계}

\begin{itemize}
  \item N\ensuremath{\approx}100의 소표본 일반화 제한
  \item 단일 대학교 편의 표집
  \item 자기보고 측정의 사회적 바람직성 편향 가능성
\end{itemize}


\subsection{6. 결론}

[데이터 수집 및 분석 완료 후 작성]


\subsection{참고문헌}

Ajzen, I., \& Fishbein, M. (1977). Attitude-behavior relations: A theoretical analysis and review of empirical research. \textit{Psychological Bulletin}, \textit{84}(5), 888–918.

Bernardo, J. M., \& Smith, A. F. M. (1994). \textit{Bayesian theory}. Wiley.

Carpenter, B., Gelman, A., Hoffman, M. D., Lee, D., Goodrich, B., Betancourt, M., \& Riddell, A. (2017). Stan: A probabilistic programming language. \textit{Journal of Statistical Software}, \textit{76}(1), 1–32.

Deci, E. L., \& Ryan, R. M. (1985). \textit{Intrinsic motivation and self-determination in human behavior}. Plenum.

교육부 (2023). \textit{2023년 교육기본통계}. 교육부.

Kaplan, R. B. (1966). Cultural thought patterns in inter-cultural education. \textit{Language Learning}, \textit{16}(1–2), 1–20.

Krashen, S. D. (1982). \textit{Principles and practice in second language acquisition}. Pergamon Press.

Masters, G. N. (1982). A Rasch model for partial credit scoring. \textit{Psychometrika}, \textit{47}(2), 149–174.

주월랑 (2022). 학문 목적 한국어 학습자의 한국어 쓰기 태도 연구. \textit{반교어문학회}, \textit{63}, 185–208.


\subsection{부록 A. 소프트웨어 및 재현 가능성}

\begin{table}[htbp]
\centering
\small
\begin{tabular}{lr}
\toprule
파일 & 역할 \\
\midrule
\texttt{t2\\_model\\_check.py} & 사전 예측 점검 및 보정값 산출 \\
\texttt{t2\\_analysis.py} & MCMC 본 분석 \\
\texttt{t2\\_figures.py} & 그림 생성 \\
\texttt{t2\\_data.csv} & 실측 설문 데이터 (연구자 수집) \\
\texttt{../sem\\_pcm\\_v2.stan} & 약한 사전 Stan 모형 \\
\texttt{../sem\\_pcm\\_with\\_prior.stan} & 강한 사전 Stan 모형 \\
\bottomrule
\end{tabular}
\end{table}

\subsection{부록 B. 설문지}

주월랑(2022) 원논문의 21개 문항을 그대로 사용. (원논문 부록 참조)

\textbf{문항 구성}

\begin{itemize}
  \item 쓰기인식(y1–y4, 4문항): 한국어 쓰기의 가치와 필요성에 대한 인식
  \item 쓰기반응(y5–y15, 11문항): 한국어 쓰기에 대한 정서적 반응 (즐거움, 자신감, 불안 등)
  \item 수행태도(y16–y21, 6문항): 실제 쓰기 수행 시의 전략적 태도 (독자 고려, 수정 등)
\end{itemize}

\end{document}
